\documentclass[11pt,a4paper]{article}

% ── FONTS ────────────────────────────────────────────────────
\usepackage{fontspec}
\setmainfont{Calibri}
\setsansfont{Calibri}

% ── LAYOUT ───────────────────────────────────────────────────
\usepackage[top=1.8cm,bottom=2cm,left=2cm,right=2cm]{geometry}
\usepackage{graphicx}
\usepackage{xcolor}
\usepackage{booktabs}
\usepackage{tabularx}
\usepackage{enumitem}
\usepackage{setspace}
\usepackage{fancyhdr}
\usepackage{titlesec}
\usepackage{hyperref}
\usepackage{microtype}
\usepackage{array}
\usepackage{colortbl}

\hypersetup{colorlinks=true,linkcolor=navyblue,urlcolor=accentteal,citecolor=navyblue}

% ── BRAND COLORS ────────────────────────────────────────────
\definecolor{navyblue}{HTML}{003566}
\definecolor{gold}{HTML}{C9A84C}
\definecolor{lightgray}{HTML}{F4F6FB}
\definecolor{darkgray}{HTML}{4A4A4A}
\definecolor{accentteal}{HTML}{0077B6}

% ── SECTION STYLING ─────────────────────────────────────────
\titleformat{\section}[block]
  {\normalfont\bfseries\fontsize{11}{13}\selectfont\color{navyblue}}
  {}{0em}{}
\titlespacing*{\section}{0pt}{10pt}{4pt}

\titleformat{\subsection}[block]
  {\normalfont\bfseries\fontsize{10}{12}\selectfont\color{accentteal}}
  {}{0em}{}
\titlespacing*{\subsection}{0pt}{6pt}{2pt}

% ── HEADER / FOOTER ─────────────────────────────────────────
\setlength{\headheight}{16pt}
\pagestyle{fancy}
\fancyhf{}
\fancyhead[L]{\footnotesize\color{navyblue}\textbf{AIGGPA Bhopal} \quad\textbar\quad
  Human Capital: Investment in Education, Skill Development and Employment}
\fancyhead[R]{\footnotesize\color{navyblue}April 2026}
\fancyfoot[C]{\footnotesize\color{navyblue}\thepage}
\renewcommand{\headrulewidth}{0.4pt}
\renewcommand{\headrule}{\color{navyblue}\hrule width\headwidth height\headrulewidth}
\renewcommand{\footrulewidth}{1pt}
\renewcommand{\footrule}{\color{gold}\hrule width\headwidth height 1pt}

% ── LISTS ───────────────────────────────────────────────────
\setlist[itemize]{leftmargin=1.2em, itemsep=1pt, topsep=2pt}

\newcommand{\rupee}{\textrm{\fontspec{Calibri}₹}}

% ────────────────────────────────────────────────────────────
\begin{document}
\onehalfspacing

\begin{center}
  {\fontsize{12}{14}\selectfont\bfseries\color{navyblue}
  Top 5 States Investing in Education, Skill Development \& Employment:\\[2pt]
  A Comparative Analysis}\\[6pt]
  {\small\color{darkgray} Aakriti \quad\textbar\quad AIGGPA Summer Intern 2026}
\end{center}

\vspace{0.2cm}

\section{1. Ranking Methodology}

States are ranked using a composite view of (a)~\textbf{absolute budget allocation} (in \rupee~Crore) for education and skill development, and (b)~\textbf{percentage of total state budget} devoted to education -- accounting for both fiscal commitment and scale of investment. Data is drawn from 2025--26 Budget Estimates and 2026--27 announcements where available.

\vspace{0.15cm}
\section{2. Top 5 States -- Budget Allocation Snapshot}

\vspace{0.1cm}
\renewcommand{\arraystretch}{1.3}
\begin{table}[h!]
\centering
\footnotesize
\caption{\textbf{Top 5 States -- Education \& Skill Development Budget (2025--26 BE)}}
\vspace{0.1cm}
\begin{tabularx}{\textwidth}{>{\columncolor{lightgray}}c l >{\centering\arraybackslash}X >{\centering\arraybackslash}X >{\centering\arraybackslash}X}
\toprule
\rowcolor{navyblue}
\textcolor{white}{\textbf{Rank}} & \textcolor{white}{\textbf{State}} & \textcolor{white}{\textbf{Edu. Budget (\rupee~Cr)}} & \textcolor{white}{\textbf{\% of State Budget}} & \textcolor{white}{\textbf{Key Strength}} \\
\midrule
1 & Uttar Pradesh      & $\sim$1,12,000  & $\sim$15.2\%  & Highest absolute allocation \\
\rowcolor{white}
2 & Maharashtra         & $\sim$1,05,000  & $\sim$14.8\%  & Strong higher education infra. \\
3 & Bihar               & $\sim$48,000    & $\sim$22.1\%  & Highest budget \% share \\
\rowcolor{white}
4 & Tamil Nadu           & $\sim$52,000    & $\sim$15.5\%  & NEP 2020 early adopter \\
5 & Karnataka            & $\sim$47,500    & $\sim$14.5\%  & IT sector--academia alignment \\
\bottomrule
\end{tabularx}
\vspace{0.05cm}
\par\scriptsize Source: Respective State Budget Documents 2025--26; PRS India State Budget Analyses; RBI State Finances Reports.
\end{table}

\vspace{0.15cm}
\section{3. State-wise Analysis}

\subsection{Uttar Pradesh}
India's most populous state leads in absolute education spending ($\sim$\rupee1.12 lakh Cr). Major programmes include \textit{Abhyudaya Coaching} for competitive exams and large-scale ITI expansion. The state recruits the highest number of primary teachers annually ($\sim$70,000 in 2024--25) and has established Skill Development Missions in all 75 districts.

\subsection{Maharashtra}
Maharashtra allocates $\sim$\rupee1.05 lakh Cr, leveraging its industrial base to drive apprenticeship-linked skilling. Pune hosts India's largest cluster of engineering colleges. The state's \textit{Mahaswayam Portal} integrates job-matching, skill training, and entrepreneurship support for youth.

\subsection{Bihar}
Despite a smaller total budget, Bihar devotes the \textbf{highest share ($\sim$22\%)} to education -- reflecting the state's prioritisation of human capital to overcome low per-capita income. The \textit{Bihar Student Credit Card Scheme} provides up to \rupee4 lakh for higher education loans at zero interest, and the state runs one of India's largest mid-day meal programmes.

\subsection{Tamil Nadu}
Tamil Nadu combines strong public university infrastructure with early NEP 2020 implementation. Its \textit{Naan Mudhalvan} (``I am the First'') platform provides free AI, blockchain, and data-science courses to college students, skilling over 15 lakh youth since launch. Education spending at $\sim$15.5\% of its budget ranks among the highest in the South.

\subsection{Karnataka}
Karnataka's Bengaluru ecosystem creates a unique industry--academia feedback loop. The state invests heavily in \textit{Karnataka Digital Economy Mission (KDEM)} and has India's highest density of engineering seats per capita. Its \textit{Skill Connect} initiative partners with 500+ employers for placement-linked training programmes.

\vspace{0.15cm}
\section{4. Key Takeaway}

Large-population states (UP, Maharashtra) dominate in absolute terms, but \textbf{proportional commitment} -- as seen in Bihar ($>$20\% of budget to education) -- is arguably a stronger indicator of policy priority. States with strong industry linkages (Karnataka, Tamil Nadu) show higher employability outcomes despite moderate spending levels.

\vspace{0.25cm}
\noindent\rule{\textwidth}{0.5pt}
{\scriptsize\textbf{References:}
(1)~PRS Legislative Research, \textit{State Budget Analyses 2025--26}, prsindia.org.
(2)~Reserve Bank of India, \textit{State Finances: A Study of Budgets 2024--25}, rbi.org.in.
(3)~Centre for Budget and Governance Accountability (CBGA), \textit{State-wise Education Expenditure}, cbgaindia.org.
(4)~National Skill Development Corporation, \textit{State Skilling Ecosystem Reports}, nsdcindia.org.
(5)~Naan Mudhalvan Portal, Government of Tamil Nadu, naanmudhalvan.tn.gov.in.
(6)~Bihar Student Credit Card Scheme, \textit{Official Portal}, 7nishchay-yuvaupmission.bihar.gov.in.
}

\end{document}
